\subsection{Release- und Packaging-Workflow}
\label{subsec:release-packaging-workflow}

\texttt{VERSION} ist die zentrale Versionsquelle. \texttt{version check}
vergleicht deren SemVer-Wert mit den Metadaten aller aktiven Frontend- und
Tauri-Komponenten; \texttt{version sync} ist der davon getrennte, schreibende
Abgleich. \texttt{release check} erzeugt und veröffentlicht selbst keine
Artefakte. Das Gate kombiniert Versionsprüfung, Produktidentität, bei Tauri
CSP und Capabilities, Tagbezug und sauberen Git-Arbeitsbaum. Bei Tauri kommt
die Anwesenheit externer Signing-Variablen hinzu. Fehlende
Signing-Konfiguration bleibt eine Warnung; inkonsistente Versionen, ein
unpassender Release-Tag, ein veränderter
Arbeitsbaum oder bekannte Template-Platzhalter im abgeleiteten Projekt sind
Fehler. Nur das Master-Repository akzeptiert seine kanonische
Template-Identität anhand des dort vorhandenen Profilmatrix-Workflows
\parencite{kleiveist2026release,kleiveist2026control}.

Der Webpfad liefert den statischen Vite-Build und das ZIP-Paket. Für Desktop
wählt ein Aufruf eine Plattformstrategie; Tauri baut auf dem jeweiligen Host
native Windows-, macOS- oder Linux-Bundles. Zusätzliche Strategien existieren
für ein portables Windows-ZIP und einen Linux-basierten Cross-Build. Die
CI-Matrix erzeugt kurzlebige, ausdrücklich unsignierte Prüfarbeitsprodukte auf
den drei nativen Betriebssystemen.

\begin{figure}[H]
  \centering
  \resizebox{0.94\textwidth}{!}{%
  \begin{tikzpicture}[node distance=7mm and 12mm]
    \node[casePrimary,minimum width=38mm] (tauri) {Tauri Build};
    \node[caseBox,minimum width=38mm,below=of tauri] (platform) {Plattform-Build};
    \node[caseBox,minimum width=28mm,below left=of platform] (windows) {Windows};
    \node[caseBox,minimum width=28mm,below=of platform] (macos) {macOS};
    \node[caseBox,minimum width=28mm,below right=of platform] (linux) {Linux};
    \node[caseBox,minimum width=38mm,below=15mm of macos] (packaging) {Packaging};
    \node[caseOptional,minimum width=72mm,below=of packaging] (external)
      {Signing / Notarisierung / Veröffentlichung\\produktspezifische oder externe Schritte};

    \draw[caseFlow] (tauri) -- (platform);
    \draw[caseFlow] (platform) -- (windows);
    \draw[caseFlow] (platform) -- (macos);
    \draw[caseFlow] (platform) -- (linux);
    \draw[caseFlow] (windows) |- (packaging);
    \draw[caseFlow] (macos) -- (packaging);
    \draw[caseFlow] (linux) |- (packaging);
    \draw[caseSecondaryFlow] (packaging) -- (external);
  \end{tikzpicture}}
  \caption{Plattformübergreifendes Modell für Desktop-Releases}
  \label{fig:desktop-release-model}
\end{figure}


Signierung, macOS-Notarisierung, Veröffentlichung, App-Stores und die
Aktivierung eines Updaters sind nicht Teil der generischen Baseline. Sie
benötigen Produktidentität, geschützte Zugangsdaten und produktspezifische
Freigaberegeln. Der Release-Workflow für Tags oder manuelle Auslösung führt
Tests, Web- und Container-Build, Release-Gate und anschließend die
unsignierte Desktopmatrix zusammen, enthält aber weder Veröffentlichungs- noch
Deployment-Job \parencite{kleiveist2026release,kleiveist2026ci}.

Für Cloudprofile kontrollieren \texttt{container doctor} und
\texttt{container validate} Docker, Compose und die Deployment-Dateien;
\texttt{build container} baut versionierte Backend- und Frontend-Images. Die
Compose-Basis dient der lokalen Produktionssimulation und hält PostgreSQL in
einem optionalen Compose-Profil. Sie legt keinen Zielanbieter fest und ersetzt
weder produktive Secret-Verwaltung noch einen kontrollierten Alembic-
Migrationsjob. Docker-/Deployment-Basis und fertige AWS-, Azure- oder
GCP-Architektur sind daher ausdrücklich nicht gleichzusetzen
\parencite{kleiveist2026container,kleiveist2026deployment}.
